\IfFileExists{stacks-project.cls}{%
\documentclass{stacks-project}
}{%
\documentclass{amsart}
}
\usepackage{amssymb}

% For dealing with references we use the comment environment
\usepackage{verbatim}
\usepackage{fancyvrb}
\newenvironment{reference}{\comment}{\endcomment}
%\newenvironment{reference}{}{}
\newenvironment{slogan}{\comment}{\endcomment}
\newenvironment{history}{\comment}{\endcomment}

% For source figures retained as assets
\usepackage{graphicx}

% For commutative diagrams we use Xy-pic
\usepackage[all]{xy}

% We use 2cell for 2-commutative diagrams.
\xyoption{2cell}
\UseAllTwocells

% We use multicol for the list of chapters between chapters
\usepackage{multicol}

% This is generally recommended for better output
\usepackage{lmodern}
\usepackage[T1]{fontenc}
\usepackage{textalpha}

% For cross-file-references
% We prefer the snapshot tmp/xref/#2.aux,
% which latexmkrc replaces atomically, so
% that documents compiled simultaneously
% never read a half-written .aux file.
% The last argument keeps links to #2.pdf.
\usepackage{xr-hyper}
\newcommand{\maybeexternaldocument}[2]{%
  \IfFileExists{#2.aux}{%
    \externaldocument[#1]{#2}%
  }{%
    \PackageWarning{stack}{Missing #2.aux; external references with prefix #1 unavailable}%
  }%
}
\let\externaldocumentorig\externaldocument
\renewcommand{\externaldocument}[2][]{%
  \IfFileExists{tmp/xref/#2.aux}{%
    \externaldocumentorig[#1]{tmp/xref/#2}[#2.pdf]%
  }{%
  \IfFileExists{#2.aux}{%
    \externaldocumentorig[#1]{#2}%
  }{%
    \PackageWarning{stack}{Missing #2.aux; external references with prefix #1 unavailable}%
  }%
  }%
}

% Package for hypertext links:
\usepackage{hyperref}

% For any local file, say "hello.tex" you want to link to please
% use \externaldocument[hello-]{hello}
\externaldocument[introduction-]{introduction}
\externaldocument[conventions-]{conventions}
\externaldocument[sets-]{sets}
\externaldocument[categories-]{categories}
\externaldocument[topology-]{topology}
\externaldocument[sheaves-]{sheaves}
\externaldocument[sites-]{sites}
\externaldocument[stacks-]{stacks}
\externaldocument[fields-]{fields}
\externaldocument[algebra-]{algebra}
\externaldocument[brauer-]{brauer}
\externaldocument[homology-]{homology}
\externaldocument[derived-]{derived}
\externaldocument[simplicial-]{simplicial}
\externaldocument[more-algebra-]{more-algebra}
\externaldocument[smoothing-]{smoothing}
\externaldocument[modules-]{modules}
\externaldocument[sites-modules-]{sites-modules}
\externaldocument[injectives-]{injectives}
\externaldocument[cohomology-]{cohomology}
\externaldocument[sites-cohomology-]{sites-cohomology}
\externaldocument[dga-]{dga}
\externaldocument[dpa-]{dpa}
\externaldocument[sdga-]{sdga}
\externaldocument[hypercovering-]{hypercovering}
\externaldocument[schemes-]{schemes}
\externaldocument[constructions-]{constructions}
\externaldocument[properties-]{properties}
\externaldocument[morphisms-]{morphisms}
\externaldocument[coherent-]{coherent}
\externaldocument[divisors-]{divisors}
\externaldocument[limits-]{limits}
\externaldocument[varieties-]{varieties}
\externaldocument[topologies-]{topologies}
\externaldocument[descent-]{descent}
\externaldocument[perfect-]{perfect}
\externaldocument[more-morphisms-]{more-morphisms}
\externaldocument[flat-]{flat}
\externaldocument[groupoids-]{groupoids}
\externaldocument[more-groupoids-]{more-groupoids}
\externaldocument[etale-]{etale}
\externaldocument[chow-]{chow}
\externaldocument[intersection-]{intersection}
\externaldocument[pic-]{pic}
\externaldocument[weil-]{weil}
\externaldocument[adequate-]{adequate}
\externaldocument[dualizing-]{dualizing}
\externaldocument[duality-]{duality}
\externaldocument[discriminant-]{discriminant}
\externaldocument[derham-]{derham}
\externaldocument[local-cohomology-]{local-cohomology}
\externaldocument[algebraization-]{algebraization}
\externaldocument[curves-]{curves}
\externaldocument[resolve-]{resolve}
\externaldocument[models-]{models}
\externaldocument[functors-]{functors}
\externaldocument[equiv-]{equiv}
\externaldocument[pione-]{pione}
\externaldocument[etale-cohomology-]{etale-cohomology}
\externaldocument[proetale-]{proetale}
\externaldocument[relative-cycles-]{relative-cycles}
\externaldocument[more-etale-]{more-etale}
\externaldocument[trace-]{trace}
\externaldocument[crystalline-]{crystalline}
\externaldocument[spaces-]{spaces}
\externaldocument[spaces-properties-]{spaces-properties}
\externaldocument[spaces-morphisms-]{spaces-morphisms}
\externaldocument[decent-spaces-]{decent-spaces}
\externaldocument[spaces-cohomology-]{spaces-cohomology}
\externaldocument[spaces-limits-]{spaces-limits}
\externaldocument[spaces-divisors-]{spaces-divisors}
\externaldocument[spaces-over-fields-]{spaces-over-fields}
\externaldocument[spaces-topologies-]{spaces-topologies}
\externaldocument[spaces-descent-]{spaces-descent}
\externaldocument[spaces-perfect-]{spaces-perfect}
\externaldocument[spaces-more-morphisms-]{spaces-more-morphisms}
\externaldocument[spaces-flat-]{spaces-flat}
\externaldocument[spaces-groupoids-]{spaces-groupoids}
\externaldocument[spaces-more-groupoids-]{spaces-more-groupoids}
\externaldocument[bootstrap-]{bootstrap}
\externaldocument[spaces-pushouts-]{spaces-pushouts}
\externaldocument[spaces-chow-]{spaces-chow}
\externaldocument[groupoids-quotients-]{groupoids-quotients}
\externaldocument[spaces-more-cohomology-]{spaces-more-cohomology}
\externaldocument[spaces-simplicial-]{spaces-simplicial}
\externaldocument[spaces-duality-]{spaces-duality}
\externaldocument[formal-spaces-]{formal-spaces}
\externaldocument[restricted-]{restricted}
\externaldocument[spaces-resolve-]{spaces-resolve}
\externaldocument[formal-defos-]{formal-defos}
\externaldocument[defos-]{defos}
\externaldocument[cotangent-]{cotangent}
\externaldocument[examples-defos-]{examples-defos}
\externaldocument[algebraic-]{algebraic}
\externaldocument[examples-stacks-]{examples-stacks}
\externaldocument[stacks-sheaves-]{stacks-sheaves}
\externaldocument[criteria-]{criteria}
\externaldocument[artin-]{artin}
\externaldocument[quot-]{quot}
\externaldocument[stacks-properties-]{stacks-properties}
\externaldocument[stacks-morphisms-]{stacks-morphisms}
\externaldocument[stacks-limits-]{stacks-limits}
\externaldocument[stacks-cohomology-]{stacks-cohomology}
\externaldocument[stacks-perfect-]{stacks-perfect}
\externaldocument[stacks-introduction-]{stacks-introduction}
\externaldocument[stacks-more-morphisms-]{stacks-more-morphisms}
\externaldocument[stacks-geometry-]{stacks-geometry}
\externaldocument[moduli-]{moduli}
\externaldocument[moduli-curves-]{moduli-curves}
\externaldocument[examples-]{examples}
\externaldocument[exercises-]{exercises}
\externaldocument[guide-]{guide}
\externaldocument[desirables-]{desirables}
\externaldocument[coding-]{coding}
\externaldocument[obsolete-]{obsolete}
\externaldocument[fdl-]{fdl}
\externaldocument[index-]{index}

% Heap Project documents
\externaldocument[k3-]{k3}
\externaldocument[torelli-]{torelli}
\externaldocument[linear-algebra-]{linear-algebra}
\externaldocument[groups-]{groups}
\externaldocument[complex-analysis-]{complex-analysis}
\externaldocument[functional-analysis-]{functional-analysis}
\externaldocument[categories-]{categories}
\externaldocument[commutative-algebra-]{commutative-algebra}
\externaldocument[ringed-spaces-]{ringed-spaces}
\externaldocument[homological-algebra-]{homological-algebra}
\externaldocument[lie-algebras-]{lie-algebras}
\externaldocument[representation-theory-]{representation-theory}
\externaldocument[smooth-manifolds-]{smooth-manifolds}
\externaldocument[differential-topology-]{differential-topology}
\externaldocument[algebraic-topology-]{algebraic-topology}
\externaldocument[differential-geometry-]{differential-geometry}
\externaldocument[algebraic-geometry-]{algebraic-geometry}
\externaldocument[symplectic-geometry-]{symplectic-geometry}
\externaldocument[complex-geometry-]{complex-geometry}
\externaldocument[hodge-theory-]{hodge-theory}
\externaldocument[%
variations-of-hodge-structures-]{%
variations-of-hodge-structures}
\externaldocument[%
homological-mirror-symmetry-]{%
homological-mirror-symmetry}
\maybeexternaldocument{deformations-}{deformations}
\externaldocument[vertex-algebras-]{vertex-algebras}
\externaldocument[springer-]{springer}
\externaldocument[infty-categories-]{infty-categories}
\externaldocument[seminars-differential-geometry-]{%
seminars-differential-geometry}
\externaldocument[seminars-complex-geometry-]{%
seminars-complex-geometry}
\externaldocument[seminars-algebraic-geometry-]{%
seminars-algebraic-geometry}
\externaldocument[seminars-others-]{%
seminars-others}
\externaldocument[%
geometric-prequantization-]{%
geometric-prequantization}
\externaldocument[surfaces-]{surfaces}
\externaldocument[birational-maps-commutative-algebra-]{birational-maps-commutative-algebra}
\externaldocument[cremona-transformations-]{cremona-transformations}
\externaldocument[derived-categories-of-sheaves-]{derived-categories-of-sheaves}
\externaldocument[geometric-invariant-theory-]{geometric-invariant-theory}
\externaldocument[geometric-stability-conditions-and-group-actions-]{geometric-stability-conditions-and-group-actions}
\externaldocument[moduli-spaces-of-sheaves-]{moduli-spaces-of-sheaves}
\externaldocument[bridgeland-stability-general-theory-]{bridgeland-stability-general-theory}
\externaldocument[hilbert-schemes-]{hilbert-schemes}
\externaldocument[noncommutative-abelian-surfaces-and-kummer-type-hyperkahler-varieties-]{noncommutative-abelian-surfaces-and-kummer-type-hyperkahler-varieties}
\externaldocument[mumford-tate-groups-in-hodge-theory-]{mumford-tate-groups-in-hodge-theory}
\externaldocument[hodge-birational-atoms-2026-]{hodge-birational-atoms-2026}
\externaldocument[xxii-egd-2026-]{%
xxii-egd-2026}
\externaldocument[pde-]{pde}
\externaldocument[probability-]{probability}

\externaldocument[linguistics-]{linguistics}
\externaldocument[genealogia-familiar-]{%
genealogia-familiar}

\externaldocument[%
16th-alga-meeting-2026-]{%
16th-alga-meeting-2026}

\externaldocument[%
iv-jornadas-lefschetz-2026-]{%
iv-jornadas-lefschetz-2026}

\externaldocument[reading-the-masters-]{%
reading-the-masters}

\externaldocument[real-analysis-]{%
real-analysis}

\externaldocument[optimal-transport-]{%
optimal-transport}

\externaldocument[history-]{history}

\externaldocument[seiberg-witten-]{seiberg-witten}

\externaldocument[history-philosophy-biology-]{history-philosophy-biology}

\externaldocument[chemistry-]{chemistry}

\externaldocument[molecular-cell-biology-]{molecular-cell-biology}

\externaldocument[microbiology-]{microbiology}

\externaldocument[protistology-]{protistology}

\externaldocument[mycology-]{mycology}

\externaldocument[botany-]{botany}

\externaldocument[zoology-]{zoology}

\externaldocument[animal-histology-]{animal-histology}

\externaldocument[animal-physiology-]{animal-physiology}

\externaldocument[systematics-biogeography-]{systematics-biogeography}

\externaldocument[ecology-natural-resources-]{ecology-natural-resources}

\externaldocument[earth-sciences-]{earth-sciences}

\externaldocument[biostatistics-]{biostatistics}

\externaldocument[mathematical-systems-biology-]{mathematical-systems-biology}

\externaldocument[cimpa-el-salvador-2024-]{cimpa-el-salvador-2024}

% Theorem environments.
%
\theoremstyle{plain}
\newtheorem{theorem}[subsection]{Theorem}
\newtheorem{proposition}[subsection]{Proposition}
\newtheorem{lemma}[subsection]{Lemma}
\newtheorem{corollary}[subsection]{Corollary}

\theoremstyle{definition}
\newtheorem{definition}[subsection]{Definition}
\newtheorem{example}[subsection]{Example}
\newtheorem{exercise}[subsection]{Exercise}
\newtheorem{situation}[subsection]{Situation}

\theoremstyle{remark}
\newtheorem{remark}[subsection]{Remark}
\newtheorem{remarks}[subsection]{Remarks}

\numberwithin{equation}{subsection}

% Macros
%
\def\lim{\mathop{\mathrm{lim}}\nolimits}
\def\colim{\mathop{\mathrm{colim}}\nolimits}
\def\Spec{\mathop{\mathrm{Spec}}}
\def\Hom{\mathop{\mathrm{Hom}}\nolimits}
\def\Ext{\mathop{\mathrm{Ext}}\nolimits}
\def\SheafHom{\mathop{\mathcal{H}\!\mathit{om}}\nolimits}
\def\SheafExt{\mathop{\mathcal{E}\!\mathit{xt}}\nolimits}
\def\Sch{\mathit{Sch}}
\def\Mor{\mathop{\mathrm{Mor}}\nolimits}
\def\Ob{\mathop{\mathrm{Ob}}\nolimits}
\def\Sh{\mathop{\mathit{Sh}}\nolimits}
\def\NL{\mathop{N\!L}\nolimits}
\def\CH{\mathop{\mathrm{CH}}\nolimits}
\def\proetale{{pro\text{-}\acute{e}tale}}
\def\etale{{\acute{e}tale}}
\def\QCoh{\mathit{QCoh}}
\def\Ker{\mathop{\mathrm{Ker}}}
\def\Im{\mathop{\mathrm{Im}}}
\def\Coker{\mathop{\mathrm{Coker}}}
\def\Coim{\mathop{\mathrm{Coim}}}

% Boxtimes
%
\DeclareMathSymbol{\boxtimes}{\mathbin}{AMSa}{"02}

%
% Macros for moduli stacks/spaces
%
\def\QCohstack{\mathcal{QC}\!\mathit{oh}}
\def\Cohstack{\mathcal{C}\!\mathit{oh}}
\def\Spacesstack{\mathcal{S}\!\mathit{paces}}
\def\Quotfunctor{\mathrm{Quot}}
\def\Hilbfunctor{\mathrm{Hilb}}
\def\Curvesstack{\mathcal{C}\!\mathit{urves}}
\def\Polarizedstack{\mathcal{P}\!\mathit{olarized}}
\def\Complexesstack{\mathcal{C}\!\mathit{omplexes}}
% \Pic is the operator that assigns to X its picard group, usage \Pic(X)
% \Picardstack_{X/B} denotes the Picard stack of X over B
% \Picardfunctor_{X/B} denotes the Picard functor of X over B
\def\Pic{\mathop{\mathrm{Pic}}\nolimits}
\def\Picardstack{\mathcal{P}\!\mathit{ic}}
\def\Picardfunctor{\mathrm{Pic}}
\def\Deformationcategory{\mathcal{D}\!\mathit{ef}}


\begin{document}

\title{Chemistry}
\maketitle

\phantomsection
\label{section-phantom}

\tableofcontents

\section{What this chapter covers}
\label{section-chem-overview}

The chemistry a biologist needs: water and
the forces between molecules, solutions and
energy, acids, bases and redox, then the
organic chemistry of carbon compounds,
isomerism and the four families of
biomolecules, and the laboratory techniques
for separating and purifying natural
products.

\textbf{Sources of these notes.}
Two courses of the Facultad de Ciencias,
UNAM. Química (1st semester, 2014-1): a few
of Daniel's laboratory discussions
(crystallization and sublimation, flame
tests, dissolution temperatures), his summary
of thermodynamics, and photographs of the
laboratory manual and of the practices on
titration and chlorine. Química Orgánica (2nd
semester, 2014-2): the official syllabus, his
answers to four problem sets, twelve
laboratory reports and a team project on
extracting pelargonidin 3-glucoside from
strawberries and menthol from mint. The
material was checked and completed with Kotz,
Treichel and Townsend, \emph{Chemistry and
Chemical Reactivity}
\cite{kotz2012chemistry}, Wade's
\emph{Organic Chemistry}
\cite{wade2013organic} (the course textbook),
Voet, Voet and Pratt
\cite{voet2013fundamentals} for water and
lipids, Pavia et al.\
\cite{pavia2002microscale} for laboratory
techniques, and Gould et al.\
\cite{gould2009anthocyanins} for
anthocyanins.

\section{Water, polarity and intermolecular
forces}
\label{section-chem-water}

A covalent bond between atoms of different
electronegativity is polar. A molecule is
polar if its bond dipoles do not cancel:
$\mathrm{CO_2}$ (linear) is nonpolar,
$\mathrm{H_2O}$ (bent) is polar \cite[Section
2-9]{wade2013organic}. In water the O--H
bonds make an angle of 104.5$^\circ$, because
the four $sp^3$ orbitals of oxygen point
roughly to the corners of a tetrahedron, two
holding hydrogens and two holding lone pairs
\cite[Chapter 2]{voet2013fundamentals}.

\textbf{Intermolecular forces.}
From weakest to strongest between neutral
molecules \cite[Section
2-10]{wade2013organic}, \cite[Chapter
12]{kotz2012chemistry}:
\begin{itemize}
\item London dispersion forces, between
instantaneous and induced dipoles: present in
all molecules, larger for larger and more
elongated molecules;
\item dipole-dipole forces, between permanent
dipoles;
\item hydrogen bonds: a hydrogen on N, O or F
attracted to a lone pair on another N, O or
F;
\item ion-dipole forces, which hydrate ions
in solution.
\end{itemize}
These forces explain boiling points and
solubilities; the first problem set asked to
sort twenty compounds by their strongest
interaction.

\textbf{Water.}
Each water molecule can make four hydrogen
bonds, two as donor and two as acceptor. This
network gives water its high boiling point,
heat capacity and heat of vaporization, and
makes ice (an open hydrogen-bonded lattice)
less dense than liquid water, so lakes freeze
from the top \cite[Chapter
12]{kotz2012chemistry}, \cite[Chapter
2]{voet2013fundamentals}.

\begin{definition}
\label{definition-chem-hydrophilic}
A substance is \emph{hydrophilic} if it
dissolves in water, because it is ionic or
polar and can interact with water by
ion-dipole forces or hydrogen bonds; it is
\emph{hydrophobic} if it is nonpolar and does
not. An \emph{amphipathic} (amphiphilic)
molecule has both a hydrophilic and a
hydrophobic part.
\end{definition}

``Like dissolves like'' \cite[Section
2-11]{wade2013organic}. Nonpolar molecules
cluster in water not because they attract
each other strongly but because water around
them is forced into ordered cages; clustering
frees this water and raises its entropy (the
hydrophobic effect). Amphipathic molecules
therefore form micelles or bilayers that hide
their hydrophobic parts, the basis of
membranes \cite[Chapter
2]{voet2013fundamentals}.

\section{Solutions and thermochemistry}
\label{section-chem-thermochemistry}

\textbf{The laws.}
Daniel's summary of the course: energy is
conserved, $\Delta E = q + w$; the enthalpy
change $\Delta H$ is the heat exchanged at
constant pressure; by Hess's law the enthalpy
change of a reaction is the sum of those of
any sequence of steps leading to it, so,
writing $P$ for the products and $R$ for the
reactants with their coefficients $n$,
$$
\Delta_r H^\circ = \sum_P n_P \, \Delta_f
H^\circ(P) - \sum_R n_R \, \Delta_f
H^\circ(R).
$$
The entropy of the universe increases in a
spontaneous process, and a perfect crystal at
0 K has zero entropy. At constant temperature
and pressure, spontaneity is decided by the
Gibbs energy $\Delta G = \Delta H - T\Delta
S$: negative means spontaneous. Away from
standard conditions $\Delta G = \Delta
G^\circ + RT \ln Q$, and at equilibrium
($\Delta G = 0$, $Q = K$) this gives $\Delta
G^\circ = -RT \ln K$ \cite[Chapters 5 and
19]{kotz2012chemistry}.

\textbf{Dissolving.}
Dissolving an ionic solid costs the lattice
energy and releases the hydration energy of
the ions; the enthalpy of solution is their
balance and can have either sign
\cite[Section 14.2]{kotz2012chemistry}.

\begin{example}
\label{example-chem-dissolution}
In the course's experiment, water at 22
$^\circ$C rose to 26 $^\circ$C when NaOH was
dissolved (exothermic) and fell to 19
$^\circ$C with $\mathrm{NaHCO_3}$
(endothermic). For NaOH, Kotz's table gives
$\Delta_f H^\circ = -425.9$ kJ/mol for the
solid and $-469.2$ kJ/mol for the aqueous
solution, the enthalpy of solution is
$$
-469.2 - (-425.9) = -43.3
$$
kJ/mol \cite[Section
14.2]{kotz2012chemistry}. (The value $-469.2$
kJ/mol is the enthalpy of formation of
aqueous NaOH, not its enthalpy of solution.)
\end{example}

\textbf{Crystallization and sublimation.}
Two other first-semester experiments: alum
crystallizing from a cooling saturated
solution (ordering costs entropy, but the
heat released to the surroundings makes
$\Delta G$ negative at low temperature), and
the sublimation of $p$-dichlorobenzene, the
solid of bathroom deodorizer blocks, which
passes directly from solid to vapour
(endothermic) and deposits as crystals on a
cold surface.

\section{Electrons, light and flame tests}
\label{section-chem-flame-tests}

Electrons in atoms occupy discrete energy
levels. Heat in a flame excites an outer
electron to a higher level; when it falls
back, the atom emits a photon of energy $E =
h\nu$ equal to the difference between the
levels, so each element emits characteristic
wavelengths. Sodium's strong line at 589 nm
gives a yellow flame and strontium salts a
red one \cite[Section
6.3]{kotz2012chemistry}.

In the course, chlorides were burned in a
methanol flame (blue on its own). Daniel
recorded potassium as yellow, strontium as
red-pink-orange, lithium as
blue-purple-orange, sodium as pink and orange
and copper(II) as green. The standard colours
are potassium pale violet, easily hidden by
traces of sodium's yellow, lithium crimson
and copper green (not checked against the
textbooks). The test identifies metals whose
excited states emit in the visible range,
mostly alkali and alkaline-earth metals and
copper.

\section{Acids, bases and titration}
\label{section-chem-acids-bases}

A Brønsted acid donates a proton and a base
accepts one; a Lewis acid accepts an electron
pair. The strength of an acid HA in water is
measured by $K_a$, or $\mathrm{p}K_a = -\log
K_a$ (smaller $\mathrm{p}K_a$, stronger acid)
\cite[Chapter 17]{kotz2012chemistry}.

In a titration a solution of known
concentration is added until it has reacted
exactly with the analyte (the equivalence
point), detected by an indicator.
Phenolphthalein, a weak acid, is colourless
in acid and pink in base, and suits
titrations of weak acids by strong bases,
whose equivalence point lies above pH 7
\cite[Chapter 18]{kotz2012chemistry}. In the
course, NaOH was first standardized against
weighed potassium hydrogen phthalate (about
0.2 g, roughly 1 mmol of a monoprotic acid),
and then used to titrate the acetic acid of
vinegar:
$$
\mathrm{CH_3COOH} + \mathrm{OH^-} \rightarrow
\mathrm{CH_3COO^-} + \mathrm{H_2O}.
$$

\section{Redox and the chemistry of chlorine}
\label{section-chem-chlorine}

In a redox reaction one species is oxidized
(loses electrons) and another reduced (gains
them). A standard reduction potential
$E^\circ$ measures the tendency of a
half-reaction to proceed as a reduction; a
reaction combining two half-reactions is
spontaneous under standard conditions if the
cell potential $E^\circ$, the potential of
the reduction minus that of the oxidation, is
positive, since $\Delta G^\circ =
-nFE^\circ$. Away from standard conditions
the Nernst equation gives $E = E^\circ -
(RT/nF)\ln Q$ \cite[Sections
20.4--20.6]{kotz2012chemistry}.

The first-semester chlorine practice used two
couples, with standard potentials 1.36 V and
1.63 V respectively \cite[Appendix
M]{kotz2012chemistry}:
$$
\begin{aligned}
\mathrm{Cl_2} + 2e^- &\rightarrow
2\,\mathrm{Cl^-},\\
2\,\mathrm{HClO} + 2\mathrm{H^+} + 2e^-
&\rightarrow \mathrm{Cl_2} +
2\,\mathrm{H_2O}.
\end{aligned}
$$
Since $1.63 > 1.36$, in acid hypochlorous
acid oxidizes chloride to chlorine; this is
how chlorine was made in a syringe (Mattson's
method) from household bleach (NaOCl) and 6 M
HCl:
$$
\mathrm{HClO} + \mathrm{H^+} + \mathrm{Cl^-}
\rightarrow \mathrm{Cl_2} + \mathrm{H_2O}.
$$
In water and above all in base the reverse
happens: chlorine disproportionates,
$$
\mathrm{Cl_2} + \mathrm{H_2O}
\rightleftharpoons \mathrm{HClO} +
\mathrm{H^+} + \mathrm{Cl^-}, \qquad
\mathrm{Cl_2} + 2\,\mathrm{OH^-} \rightarrow
\mathrm{ClO^-} + \mathrm{Cl^-} +
\mathrm{H_2O}.
$$
Bubbling chlorine into water lowered the pH,
and silver nitrate then precipitated AgCl,
showing the chloride. Chlorine bleached fruit
juices (by oxidizing their pigments) and
oxidized bromide and iodide to
$\mathrm{Br_2}$ and $\mathrm{I_2}$, as
expected from the oxidizing strength of the
halogens, $\mathrm{F_2} > \mathrm{Cl_2} >
\mathrm{Br_2} > \mathrm{I_2}$ ($E^\circ$ =
2.87, 1.36, 1.08 and 0.535 V) \cite[Section
21.10]{kotz2012chemistry}.

\section{Carbon, hybridization and functional
groups}
\label{section-chem-carbon}

Carbon forms four bonds. Mixing its 2s
orbital with two or three 2p orbitals gives
hybrid orbitals whose number equals the
number of $\sigma$ bonds and lone pairs
around the atom \cite[Sections 2-4 and
2-6]{wade2013organic}:

\begin{center}
\begin{tabular}{llll}
Hybrid & Geometry & Angle & Example \\
\hline
$sp^3$ & tetrahedral & 109.5$^\circ$ & ethane
\\
$sp^2$ & trigonal & 120$^\circ$ & ethylene \\
$sp$ & linear & 180$^\circ$ & acetylene \\
\end{tabular}
\end{center}

A $\sigma$ bond has its electron density
along the axis between two atoms (not two
molecules, as one answer in the problem sets
had it); a $\pi$ bond is formed by
side-by-side overlap of parallel p orbitals
above and below that axis. A double bond is
one $\sigma$ and one $\pi$ bond, a triple
bond one $\sigma$ and two $\pi$; $\pi$ bonds
prevent rotation \cite[Sections 2-2 and
2-3]{wade2013organic}.

\textbf{Families.}
The syllabus covered IUPAC names and physical
properties of \cite[Sections 2-12 to
2-14]{wade2013organic}:

\begin{center}
\begin{tabular}{ll}
Family & Group \\
\hline
alkane, alkene, & C--C, C=C, \\
alkyne & C$\equiv$C \\
aromatic & benzene ring \\
alcohol, phenol & --OH \\
ether & C--O--C \\
thiol, sulfide & --SH, C--S--C \\
amine & --NH$_2$, --NHR, --NR$_2$ \\
aldehyde & --CHO \\
ketone & C--CO--C \\
carboxylic acid & --COOH \\
ester & --COOR \\
amide & --CONH$_2$ \\
\end{tabular}
\end{center}

Conjugation (alternating single and double
bonds) lets $\pi$ electrons delocalize. A
molecule with delocalized electrons is
described by several resonance structures,
none of which exists by itself; the real
molecule, the resonance hybrid, is more
stable than any of them \cite[Section
1-9]{wade2013organic}.

\section{Acidity and basicity of organic
molecules}
\label{section-chem-organic-acidity}

The more stable the conjugate base, the
stronger the acid. Three effects were studied
\cite[Sections 1-13 and
20-4]{wade2013organic}:
\begin{itemize}
\item resonance: acetic acid ($\mathrm{p}K_a$
4.74) is far more acidic than ethanol because
the negative charge of acetate is shared by
two oxygens;
\item induction: electronegative substituents
pull electron density through the $\sigma$
bonds and stabilize the anion, so
chloroacetic acid ($\mathrm{p}K_a$ 2.86) is
stronger than acetic acid, and the effect
fades with distance;
\item steric effects: bulky groups hinder the
solvation of an ion and so change acidity and
basicity.
\end{itemize}
Amines are the common organic bases; amino
acids carry both an acidic and a basic group.

\section{Isomerism and chirality}
\label{section-chem-isomerism}

\emph{Isomers} have the same molecular
formula. \emph{Constitutional} isomers differ
in which atoms are bonded to which.
\emph{Stereoisomers} have the same
connectivity and differ in the arrangement in
space: \emph{conformations} interconvert by
rotation about single bonds (staggered and
eclipsed ethane, chair cyclohexane) and are
usually not isolable; \emph{configurations}
change only by breaking bonds \cite[Sections
2-8 and 3-7]{wade2013organic}.

\begin{definition}
\label{definition-chem-chirality}
An object is \emph{chiral} if it is not
superimposable on its mirror image.
\emph{Enantiomers} are stereoisomers that are
mirror images of each other;
\emph{diastereomers} are stereoisomers that
are not. A \emph{racemic mixture} contains
equal amounts of two enantiomers. A
\emph{meso compound} has asymmetric carbons
but is achiral, because it has an internal
mirror plane.
\end{definition}

An asymmetric carbon (bearing four different
groups) is the usual cause of chirality.
Enantiomers have identical physical
properties except the direction in which they
rotate plane-polarized light: dextrorotatory
(+) or levorotatory ($-$); a racemate is
optically inactive. Diastereomers differ in
all physical properties \cite[Chapter
5]{wade2013organic}.

\textbf{Descriptors.}
\begin{itemize}
\item R/S (Cahn-Ingold-Prelog): rank the four
groups by atomic number at the first point of
difference; with the lowest-ranked group
pointing away, the sequence 1, 2, 3 is
clockwise for R and counterclockwise for S
\cite[Section 5-3]{wade2013organic}. The sign
of rotation cannot be deduced from R or S.
\item D/L: relative configuration by
comparison with glyceraldehyde in the Fischer
projection; in a sugar the asymmetric carbon
farthest from the carbonyl has its OH on the
right in the D series. Natural sugars are D,
and natural amino acids are L \cite[Sections
5-10 and 23-3]{wade2013organic}.
\item E/Z: for a double bond, rank the two
groups on each carbon by the same rules; Z
when the higher-ranked groups are on the same
side, E when opposite. It replaces cis/trans
when the latter is ambiguous \cite[Section
7-5]{wade2013organic}.
\end{itemize}
Cis-trans isomerism also occurs in rings, as
in steroids and cyclic terpenes.

\section{Biomolecules}
\label{section-chem-biomolecules}

\textbf{Amino acids and proteins.}
An $\alpha$-amino acid has an amino group and
a carboxyl group on the same carbon, with a
side chain that defines it; all but glycine
are chiral, and natural ones are L. In water
near neutral pH it is a \emph{zwitterion}
($\mathrm{H_3N^+}$--CHR--$\mathrm{COO^-}$).
The \emph{isoelectric point} pI is the pH at
which the average net charge is zero; for an
amino acid without ionizable side chain it is
the mean of the two $\mathrm{p}K_a$ values.
In electrophoresis an amino acid moves
towards the cathode below its pI and towards
the anode above it \cite[Sections 24-2 to
24-4]{wade2013organic}. The peptide bond is
an amide; resonance gives it partial
double-bond character, so it is planar and
does not rotate. The primary structure is the
sequence of residues, written from the
N-terminus \cite[Section
24-8]{wade2013organic}. Disulfide bridges
form by oxidation of two cysteine thiols.
Proteins are classed as simple or conjugated
(with a non-amino-acid part: a cofactor,
which may be a metal ion or an organic
coenzyme; a tightly or covalently bound one
is a prosthetic group), and as fibrous
(structural, elongated) or globular.

\textbf{Carbohydrates.}
Monosaccharides are polyhydroxy aldehydes
(aldoses) or ketones (ketoses). An alcohol
adds to an aldehyde to give a hemiacetal, and
to a ketone to give a hemiketal; a second
alcohol then gives an acetal (from an
aldehyde) or a ketal (from a ketone; IUPAC
calls both acetals) \cite[Section
18-17]{wade2013organic}. A sugar's own OH
group (at C5 for glucose) attacks its
carbonyl and closes a six-membered pyranose
or five-membered furanose ring; the former
carbonyl carbon becomes a new asymmetric
centre, the \emph{anomeric carbon}, giving
the $\alpha$ and $\beta$ anomers (in
$\beta$-D-glucopyranose all OH groups are
equatorial). Rings are drawn as flat Haworth
projections or, more realistically, as chairs
\cite[Sections 23-6 and
23-7]{wade2013organic}. Replacing the
anomeric OH by an OR group gives a glycoside,
an acetal; glycosidic bonds link
monosaccharides into disaccharides and
polysaccharides \cite[Section
23-11]{wade2013organic}.

\textbf{Nucleic acids.}
A nucleoside is a base joined to ribose or
deoxyribose by an N-glycosidic bond at the
anomeric carbon; a nucleotide is its
5'-phosphate; nucleotides are joined by
phosphodiester bonds. In DNA, A pairs with T
by two hydrogen bonds and G with C by three
\cite[Sections 23-19 to
23-21]{wade2013organic}. Two problem set
exercises: the strand complementary to
5'-GGCTAATCCGT-3' is 5'-ACGGATTAGCC-3'; a sea
urchin DNA with 32\% A has 32\% T and 18\%
each of G and C.

\textbf{Lipids.}
Lipids are defined by solubility in nonpolar
solvents, not by structure \cite[Chapter
25]{wade2013organic}, \cite[Section
9-1]{voet2013fundamentals}:
\begin{itemize}
\item hydrolysable (with ester or amide
bonds): waxes (a long-chain acid esterified
to a long-chain alcohol; spermaceti is cetyl
palmitate), triacylglycerols (fats if solid
at room temperature, oils if liquid:
saturated chains such as stearic acid pack
well, while cis double bonds, as in oleic
acid, kink the chains and lower the melting
point; saponification by NaOH gives glycerol
and soaps), phosphoglycerides (glycerol with
two fatty acids and a phosphate, often
esterified to an amino alcohol), and
sphingolipids, built on ceramide (sphingosine
with a fatty acid in amide bond):
sphingomyelins carry phosphocholine,
cerebrosides a single sugar and gangliosides
an oligosaccharide with sialic acid;
\item nonhydrolysable: steroids (four fused
rings, as in cholesterol), terpenes (built of
five-carbon isoprene units; monoterpenes have
ten carbons), prostaglandins (from a
20-carbon fatty acid, with a five-membered
ring) and fat-soluble vitamins.
\end{itemize}

\section{Separation and purification
techniques}
\label{section-chem-lab-techniques}

The organic chemistry laboratory taught the
classical methods, which Pavia et al.\
describe in detail
\cite{pavia2002microscale}.

\textbf{Miscibility and solubility.}
Hexane was mixed with ethyl acetate,
petroleum ether (a mixture of alkanes, not an
ether \cite[Chapter
10]{pavia2002microscale}), chloroform and
methanol: all mixed except methanol. Nine
solutes (tartaric, tannic, pyruvic and
salicylic acids, urea and plant and cochineal
material) were tested in four solvents: of 36
combinations, 11 dissolved, 11 partly and 14
not. Solubility follows polarity and hydrogen
bonding.

\textbf{Crystallization.}
The ideal solvent dissolves the compound when
hot but not when cold, while impurities
either stay dissolved on cooling or do not
dissolve at all and are filtered off hot;
slow cooling gives purer crystals
\cite[Chapter 11]{pavia2002microscale}.

\textbf{Distillation.}
Simple distillation separates liquids well
only when their boiling points differ by a
lot (more than about 100 $^\circ$C) or when
one is a minor impurity; otherwise fractional
distillation, with a column in which vapour
condenses and re-evaporates many times, is
needed. Some mixtures form azeotropes, which
distil at constant composition and
temperature and cannot be separated by
distillation \cite[Chapters 14 and
15]{pavia2002microscale}. In the course, an
unknown mixture of alcohols and one other
liquid gave a distillation curve whose
plateaus were interpreted as methanol (65
$^\circ$C), ethanol (78 $^\circ$C) and
2-methylpropan-2-ol (83 $^\circ$C), and a
final rise to 91 $^\circ$C as water, which
could not be boiled off before the heating
had to stop. Steam distillation codistils
water-immiscible compounds below 100
$^\circ$C: the mixture boils when the sum of
the two vapour pressures reaches atmospheric
pressure (Dalton's law), which protects
heat-sensitive essential oils \cite[Chapter
18]{pavia2002microscale}. Distillation at
reduced pressure lowers boiling points
\cite[Chapter 16]{pavia2002microscale}; the
rotary evaporator uses it to remove solvents
gently.

\textbf{Extraction.}
A solute partitions between two immiscible
solvents according to its relative
solubility; shaking in a separatory funnel
and repeating with several small portions
extracts it efficiently \cite[Chapter
12]{pavia2002microscale}. Maceration is
solid-liquid extraction by prolonged soaking.

\textbf{Chromatography.}
Components partition between a stationary
phase (silica gel, polar) and a mobile phase
(the eluent). Polar compounds adsorb strongly
and move little; a more polar eluent moves
everything further. In thin-layer
chromatography the retention factor $R_f$,
the distance travelled by a compound divided
by the distance travelled by the solvent
front, is characteristic of a compound only
for a given adsorbent, solvent, layer
thickness and sample load \cite[Chapter
20]{pavia2002microscale}. Colourless spots
are revealed under ultraviolet light or with
reagents; the Liebermann-Burchard reagent
(acetic anhydride and sulfuric acid) was used
in the course to reveal terpenes. Column
chromatography applies the same principle on
a preparative scale and collects fractions
\cite[Chapter 19]{pavia2002microscale}. In a
first TLC practice, a hexane extract of green
leaves separated on silica with hexane-ethyl
acetate 4:1 into several green bands
(chlorophylls) and a yellow-orange band near
the front, the least polar pigment
(carotenoids); with pure hexane only the
yellow band moved.

\section{The strawberry and mint project}
\label{section-chem-project}

The team project of the organic chemistry
course aimed to extract and purify two
natural products: the main anthocyanin of
strawberry and the main component of mint
oil.

\textbf{Anthocyanins.}
Anthocyanins are flavonoids: a 15-carbon
skeleton of two aromatic rings joined by a
three-carbon unit closed into an
oxygen-containing ring (C6-C3-C6),
glycosylated on hydroxyl groups. The coloured
form usually drawn is the flavylium cation,
stable in strong acid; at the mildly acidic
pH of vacuoles (pH 3--6) it is easily
converted into quinonoidal bases, water
adducts and chalcones, which changes or
removes the colour \cite[Chapter
3]{gould2009anthocyanins}. In strawberry the
main anthocyanin is pelargonidin
3-O-glucoside, followed by cyanidin
3-O-glucoside \cite[Section
5.4]{gould2009anthocyanins}; the two differ
by one hydroxyl group on the B ring (3' in
cyanidin).

The procedure: 1 kg of strawberries was
ground and macerated for seven days in 300 mL
of methanol with 1\% HCl (acid to keep the
flavylium form), filtered (sieve, gauze, then
vacuum; about half the material was lost),
concentrated in a rotary evaporator and
freeze-dried, giving 4.65 g of crude pigment
mixture. TLC on silica with butanol-acetic
acid-water 4:1:5 gave spots at $R_f$ 0.97,
0.76 and 0.35 with a continuous streak
between the last two, against a literature
value of 0.44 for pelargonidin 3-glucoside in
the same eluent. A silica column with the
same eluent ran extremely slowly and gave
only two fractions; the more coloured, lower
fraction presumably held the pigment, but
pure pelargonidin 3-glucoside was not
isolated.

\textbf{Menthol.}
Menthol is a monocyclic monoterpene alcohol,
$\mathrm{C_{10}H_{20}O}$, with three
asymmetric carbons and hence $2^3 = 8$
stereoisomers; natural mint oil is rich in
($-$)-menthol. From 200 g of mint leaves, the
essential oil was obtained by steam
distillation, separated from water by
liquid-liquid extraction into hexane, and
concentrated by fractional distillation. TLC
on silica with toluene-acetic acid 9:1,
revealed with the Liebermann-Burchard
reagent, gave one main spot at $R_f$ 0.67,
far from the literature value of 0.28 for
menthol in that system. Possible reasons,
discussed in the report: a different silica,
overloaded or impure samples. The mint smell
made the presence of menthol likely, but
neither product was identified with
certainty.

\bibliography{my}
\bibliographystyle{amsalpha}

\end{document}
